% Setup
\section{Setup}\label{sec:setup}

\paragraph{Assignment.}
The team was assigned the \TechNode{} technology (model file \ModelFile{}, a PTM card for the
BSIM4 model, \texttt{level=54}), a nominal supply of \VDDnom{} and the \TestCircuit{} as test
circuit. Every transistor uses the minimum channel length of the node, $L = L_{\min} = \Lmin$.
All of these values are read from \texttt{configs/project.yaml} and never hard-coded, so the
same code would run unchanged for another assignment.

\paragraph{Simulator.}
All SPICE results come from \SimVersion{} (BSIM4 \texttt{level=54}), one DC operating point
(\texttt{op}) per deck. Every deck, from a single transistor to the full test circuit, contains
the options line required by the brief,
\texttt{.options gmin=1e-18 abstol=1e-18 reltol=1e-6 vntol=1e-9}, and prints with
\texttt{set numdgt=12}, because ngspice's default precision is too coarse for currents that span
more than ten decades.

\paragraph{Model-card override (\texttt{rgatemod=0}).}
The card ships with \texttt{rgatemod=1}, which adds an internal gate-resistance node. We set
\texttt{rgatemod=0} for every deck. At DC no current flows through that resistance apart from
the tunnelling current itself, so the selector does not change the physics of leakage, but with
it enabled 15 of the 3\,600 AND2/OR2 gate decks failed to converge and ngspice reported $I_G = 0$
for many single-transistor rows. With \texttt{rgatemod=0} every deck converges, and the gate
leakage moves by only about $4\times10^{-6}$ (relative, median), far below any error discussed in
this report.
% Source of the 15/3600 and 4e-6 figures: the investigation documented in src/leakml/pdk.py
% (MODEL_OVERRIDES). They come from a one-off comparison, not from a pipeline output.

\paragraph{Gate sizing.}
Every NMOS in every gate and in the test circuit has $W = \NmosWMult\,L_{\min}$ and every PMOS
$W = \PmosWMult\,L_{\min}$. The PMOS is twice as wide because hole mobility is roughly half the
electron mobility (nominal \texttt{u0} of the card), which gives roughly balanced pull-up and
pull-down strength. Both widths lie inside the sampled range $W \in
[\WMultMin, \WMultMax]\,L_{\min}$, so the transistor models never extrapolate in $W$ at gate level.

\paragraph{Software and reproduction.}
The pipeline is Python \PythonVersion{} with NumPy \VerNumpy, pandas \VerPandas,
scikit-learn \VerScikitLearn{} and LightGBM \VerLightgbm{}; the exact environment of the final run
is recorded in \texttt{artifacts/env.json}. A single \texttt{make all} regenerates every
deliverable from the model file by running fourteen scripts in order:
environment check and extraction of the nominal BSIM4 parameters (00--01), transistor data and
data check (02--03), model comparison, final models and learning curve (04--05b), the shared PVT
corners (06), gate SPICE, gate ML and gate report (07--09), and the same three steps for the test
circuit (10--12). All randomness is seeded from \texttt{configs/project.yaml}: transistor data
\SeedTransistorData, data split \SeedSplit, model training \SeedTraining{} and PVT corners
\SeedCorners. Each transistor case draws from its own child stream of the data seed, so
regenerating one case leaves the other three bit-identical.
